% ============================================================
% Section 3 — Literature Review
% FENG 498 Graduation Project
% Izmir University of Economics
% ============================================================

\section{Literature Review}\label{sec:lit}

We searched the literature along four threads: storage location assignment
methods (especially class-based and turnover-driven approaches), ABC and
multi-criteria classification schemes, discrete-event simulation (DES)
applied to warehousing, and kitting together with in-plant logistics.
For each study we describe what was done, what KPIs were measured, and
where the scope or method differs from the present work.
The section closes with the research gap that motivated our line-aware
slotting policy.

% ------------------------------------------------------------
\subsection{Storage Location Assignment}\label{subsec:sla}
% ------------------------------------------------------------

The question of \emph{where} to put each stock-keeping unit in a warehouse
has been studied for more than six decades, and the core trade-offs were
already visible in the earliest analytical work.

\medskip
\noindent\textbf{Class-based versus dedicated versus random assignment.}
Hausman, Schwarz, and Graves~\cite{hausman1976} analyzed an automated
warehouse and compared three pure strategies: dedicated (each item assigned
one fixed location), random (locations chosen uniformly at random), and
class-based (items partitioned into classes by turnover, each class
occupying a zone).
Their key result is that a three-class policy captures most of the travel
savings of full dedicated assignment while requiring far less location space
and no re-slotting as demand mix shifts.
We work in a manual rack warehouse rather than an automated system, but the
same three-way taxonomy organizes our six policy variants.
The Baseline (Actual~SAP) policy is a historically evolved dedicated
assignment; the Heuristic policy is a rule-of-thumb dedicated assignment;
Usage-based~ABC and Double~ABC are class-based; and Line-aware Slotting
is a structured dedicated assignment that replaces the turnover criterion
with a line-proximity criterion.

\medskip
\noindent\textbf{The cube-per-order index.}
Heskett~\cite{heskett1963} proposed ordering items by the ratio of storage
volume to order frequency — the cube-per-order index (COI) — and placing
low-COI items nearest the input/output point.
Items that move often in small packages should occupy the slots with the
shortest travel path.
Our Travel-distance Optimized policy is the direct descendant of COI: it
sorts materials by pick frequency and assigns the most-picked materials to
the lowest-travel-cost bins.
As we discuss in Section~\ref{sec:results}, this strategy \emph{degrades}
lead time in our setting because assigning all high-frequency materials to
the same zone concentrates reach-truck demand spatially and creates a
queueing bottleneck that the COI analysis ignores.

\medskip
\noindent\textbf{Interactions between storage and routing policy.}
Petersen and Aase~\cite{petersen2004} used simulation to evaluate how
storage assignment and travel-routing (traversal, return, combined,
largest-gap, midpoint) interact in a manual warehouse.
They found that the best combination depends heavily on aisle structure and
that no assignment policy dominates across all routing algorithms.
Our warehouse has a different topology — a main kit corridor between
production-line-facing rack faces rather than a classical single-block
layout with dedicated pick aisles — but the lesson holds: the slotting
decision cannot be decoupled from the physical direction of kit travel
(East-facing versus West-facing bays, hand-reachable levels versus
reach-truck-only upper levels).

\medskip
\noindent\textbf{Practical principles: golden zone and fast-mover proximity.}
Frazelle~\cite{frazelle2002} and Bartholdi and Hackman~\cite{bartholdi2014}
both document the golden-zone heuristic: the body-height band of a rack
(roughly levels~3 through~5) should hold the fastest-moving items because
it minimizes operator bending and eliminates the need for powered equipment.
Bartholdi and Hackman further note that demand-weighted travel distance,
not mere frequency rank, governs the savings.
Our Line-aware Slotting policy puts this directly into practice: items
consumed by nearby production lines go into levels~2 through~5 of the
closest rack bay so that a standing operator can pick without calling a
reach truck.
The simulation results confirm the golden-zone logic — reach-truck
utilization under Line-aware Slotting falls to about \SI{1.3}{\percent},
compared to \SI{21.7}{\percent} under the current SAP baseline.

\medskip
\noindent\textbf{Travel time as the dominant cost.}
De Koster, Le-Duc, and Roodbergen~\cite{dekoster2007} surveyed order-picking
literature and reported that travel accounts for roughly half to two-thirds
of total picking time in most manual warehouses.
Our video time-motion study (the F400 study) found similar proportions for
this facility and supplied the per-activity timing distributions we used in
the simulation model.
De Koster et al.\ also identify shared-resource contention as an
under-studied cost driver in warehouse models — which is exactly the gap we
address by modelling reach trucks and operators as finite, competing
resources.

% ------------------------------------------------------------
\subsection{ABC Classification and Multi-Criteria Variants}\label{subsec:abc}
% ------------------------------------------------------------

ABC classification assigns items to priority classes — A, B, or C —
based on a single ranked criterion, most commonly annual consumption value
following the Pareto principle.
Class-A items, which are few but account for most of the throughput, go into
the most accessible locations.
In a warehouse context the access criterion shifts from value to pick
frequency or pick volume.

\medskip
\noindent\textbf{Single-criterion limits.}
Flores and Whybark~\cite{flores1986} pointed out that ranking by a single
criterion conflates items that are simultaneously high-volume \emph{and}
high-value with items that are one but not the other.
They proposed a two-dimensional matrix (a $3 \times 3$ or $5 \times 5$
grid of frequency versus value classes) to produce more differentiated
assignment priorities.
We put this idea into practice as our Double~ABC policy: materials are
ranked by an order-frequency class on one axis and a pick-volume class on
the other, and the resulting nine cells are collapsed to three priority
tiers that drive slot assignment.
In our simulation, Double~ABC yields a mean kitting lead time of about
\SI{14.7}{\minute} — essentially the same as the plain Heuristic policy
(\SI{14.7}{\minute}) and materially worse than Line-aware Slotting
(\SI{6.9}{\minute}).
We return to why multi-criteria ABC alone is insufficient in
Section~\ref{sec:results}; briefly, frequency and volume are still
aggregate measures that say nothing about which production line consumes a
given material, so they cannot prevent the cross-warehouse travel that is
the main cost in our floor plan.

\medskip
\noindent\textbf{The plant's current ABC/FMR tool.}
The Schneider Electric Manisa facility uses an SAP-based FMR (Fast/Medium/
Rare) classification that sorts materials by aggregate consumption across
all nine product families.
A component consumed heavily by Line~F400 but rarely by any other line still
receives a high FMR rank and lands in a general fast-mover zone that may be
distant from Line~F400's kitting table.
Our ZWM92 dispatch log contains about 40,800 kit orders across four months
and resolves per-line consumption at the component level, which makes it
possible to build a line-aware assignment that the plant's existing tool
cannot produce from its aggregated FMR scores.

% ------------------------------------------------------------
\subsection{Simulation-Based Warehouse Studies}\label{subsec:sim}
% ------------------------------------------------------------

Simulation has become the standard evaluation tool for warehouse design
decisions because analytical models cannot capture the stochastic
interactions among arrivals, routing, and resource sharing that dominate
real performance.

\medskip
\noindent\textbf{Scope of the literature.}
Negahban and Smith~\cite{negahban2014} reviewed simulation studies of
manufacturing systems and warehouses published over roughly two decades.
They noted a steady shift from queuing-analytic models toward DES and
agent-based models as computing costs fell, and identified resource
contention — shared equipment, shared aisles — as the most common reason
that analytic approaches under-predict congestion.
Our study sits squarely in the DES strand they describe: we used SimPy to
model arrivals as a fitted batch-arrival process, shared reach trucks and
operators as finite-capacity resources, and replicated the experiment twenty
times to obtain confidence intervals.

\medskip
\noindent\textbf{DES methodology.}
Banks et al.~\cite{banks2010} and Law~\cite{law2015} are the standard
methodology references for DES studies.
Both texts prescribe fitting input distributions from real data, using
common random numbers (CRN) across scenarios to reduce variance in
comparisons, and running enough replications for interval estimates.
We followed these prescriptions: arrival inter-batch gaps and batch sizes
are drawn from empirical distributions fitted to the ZWM92 timestamped
records, service times use lognormal distributions whose parameters come
from the F400 video study, and all six policies share the same 20~CRN
replication seeds.

\medskip
\noindent\textbf{Slotting evaluated by DES.}
Ardiansyah et al.~\cite{ardiansyah2024} applied DES in ProModel to evaluate
warehouse layout and picking efficiency improvements in an electronics
distribution centre, finding travel-time reductions of several percentage
points when items were regrouped by pick frequency.
Our study differs in three respects: we model an \emph{inbound kitting}
process rather than outbound order picking, our batched arrival stream comes
directly from four months of SAP dispatch records rather than synthetic
input, and we evaluate six policies with paired-replication $t$-tests
rather than a single before-and-after comparison.
Hapsari, Arlianto, and Santoso~\cite{hapsari2025} used DES to improve layout
and inventory control in a pharmaceutical warehouse, showing that simulation
can surface bottlenecks invisible in static ABC analysis.
We make a similar point, but our scenario adds shared mechanical handling
resources whose contention is the dominant lead-time driver.

\medskip
\noindent\textbf{Hybrid simulation and optimization.}
Amorim-Lopes et al.~\cite{amorim2021} combined probabilistic modelling,
mathematical optimization, and DES in sequence to improve picking at a large
retailer warehouse, using the DES stage to validate the optimizer's plan
under realistic congestion.
Their architecture is more elaborate than ours; we chose a pure simulation
approach because the six policies are structurally distinct enough to
evaluate by experiment, and the real SAP data provide a direct validation
target — the current SAP baseline replicates the observed throughput of
about 396~orders per active day within \SI{5}{\percent}.
Kokroko, Kobi, and Yeboah~\cite{kokroko2025} combined simulation with
optimization algorithms to maximize storage efficiency and minimize travel
distances, reporting improvements in both KPIs simultaneously.
Their objective function does not include downstream resource queueing; in
our warehouse the Travel-distance Optimized policy achieves the
third-shortest walk distance yet produces the longest lead time
(\SI{24.6}{\minute}) because it floods the reach trucks.

\medskip
\noindent\textbf{Analytical and bi-level optimization.}
Moghadam et al.~\cite{moghadam2025} proposed a bi-level optimization model
that integrates warehouse design, storage location assignment, and order
picking into a single mathematical programme.
The upper level selects storage locations; the lower level minimizes travel
on those locations.
This is a clean analytical complement to our simulation-based approach.
The limitation of such formulations is that they typically assume
deterministic or steady-state demands and do not model the stochastic
arrivals, resource contention, or shift breaks that generate the queueing
we observe.

\medskip
\noindent\textbf{Automated storage and retrieval — Kardex.}
Roodbergen and Vis~\cite{roodbergen2009} surveyed automated storage and
retrieval systems (AS/RS), including horizontal and vertical carousels of
the kind our Kardex units represent.
They report that carousel throughput is sensitive to dwell-point strategy
(where the carousel parks between requests) and that batching requests
within one cycle increases throughput substantially.
We modelled our four Kardex units as a shared resource with a single
carousel-rotation overhead per visit and multiple picks within one hold,
which is the batching strategy Roodbergen and Vis recommend.
Validation showed that Kardex utilization remains below
\SI{10}{\percent} across all policies because the approximately 2,870
Kardex-assigned materials face relatively modest aggregate demand.

% ------------------------------------------------------------
\subsection{Kitting and In-Plant Logistics}\label{subsec:kitting}
% ------------------------------------------------------------

Our warehouse does not perform outbound distribution; it feeds nine
production lines by assembling component kits and delivering them by
milkrun train.
The kitting and line-stocking literature is therefore directly relevant.

\medskip
\noindent\textbf{Kitting time efficiency.}
Hanson and Medbo~\cite{hanson2012} studied manual assembly kitting at
Swedish manufacturing plants and found that kit preparation time is
dominated by walking and searching rather than the physical act of picking,
and that grouping components by line reduces total kit assembly time
materially.
This matches what we observed: the F400 video study shows that an
operator's non-value-added travel between bays during a single kit pick is
the dominant activity.
Line-aware Slotting targets this directly by assigning a production line's
highest-demand components to the rack bays facing that line's kitting
table, so that most picks in a kit are completed without leaving a short
corridor segment.

\medskip
\noindent\textbf{Kitting versus line stocking.}
Limere et al.~\cite{limere2012} built a cost model for the choice between
kitting (components collected into a kit before delivery to the line) and
line stocking (components stored directly at the line) in automotive
assembly.
Their model shows that kitting is preferred when components are shared
across multiple products, when line-side storage space is scarce, and when
order volumes vary.
The Manisa medium-voltage switchgear warehouse operates under exactly those
conditions: components are shared across nine product families, the
production floor has no significant line-side storage, and order volumes
follow a batch-arrival pattern with high inter-batch variability — the
ZWM92 data show a coefficient of variation of about 2.0 for inter-batch
gaps within a shift.
Kitting is therefore not a design choice we evaluate here; it is the fixed
operational context within which we optimize slot assignments.

\medskip
\noindent\textbf{Milkrun scheduling.}
Emde and Boysen~\cite{emde2012} studied routing and scheduling of tow-train
(milkrun) deliveries for just-in-time supply of mixed-model assembly lines.
They showed that the sequence and timing of milkrun routes interact with
kit-assembly completion times: a milkrun that arrives before a kit is ready
either waits or makes a redundant loop.
In our model the milkrun trains are the downstream customers of the kitting
process; we model them as a resource pool whose members ferry completed kits
to the lines.
We do not optimize milkrun routes in this study — routes are fixed by the
floor plan — but the Emde-Boysen framing clarifies why lead time is the
right primary KPI: it measures the interval from kit-order release to kit
delivery at the line, which is exactly the quantity that drives line-side
buffer requirements.

% ------------------------------------------------------------
\subsection{Research Gap}\label{subsec:gap}
% ------------------------------------------------------------

The storage-assignment literature, from the early COI work of
Heskett~\cite{heskett1963} through the modern simulation studies reviewed
above, optimizes for some form of aggregate travel distance or throughput.
Even the multi-criteria ABC extension of Flores and
Whybark~\cite{flores1986} retains a single demand-weighted score per item:
it does not distinguish \emph{which demand} drives the score.
In a plant warehouse that feeds multiple distinct production lines, this
aggregation matters.
Assigning a component to the nearest-to-depot slot minimizes average travel
across all orders but may force a long detour for orders from the line at
the far end of the warehouse.

The coupling between slotting and shared handling equipment is equally
overlooked.
The Travel-distance Optimized policy in our study illustrates the problem:
by packing all high-demand items into the lowest-travel-cost zone, it
creates a reach-truck hotspot that nearly doubles lead time relative to the
SAP baseline, despite a shorter average walk.
Analytical and optimization-based methods~\cite{moghadam2025,kokroko2025}
typically assume uncapacitated or single-resource travel; the queueing
effect only appears when a DES model captures concurrent orders competing
for the same reach truck at the same moment.

So the gap is this: no study we found explicitly co-optimizes slot
assignment with (a)~the identity of the production line that will consume
each kit, and (b)~the contention for shared mechanical handling resources,
evaluated on a multi-month company dispatch log with paired statistical
testing.
Our Line-aware Slotting policy and the supporting DES model are designed
to fill that gap.
The policy assigns each material to the rack face nearest its primary
consuming line and to a manually reachable level wherever possible, so that
reach-truck demand is minimized by design.
Across twenty replications, it achieves a mean kitting lead time of about
\SI{6.9}{\minute} — about \SI{44}{\percent} below the \SI{12.4}{\minute}
recorded under the current SAP placement — while holding throughput
constant at about 399 orders per day.
